% !PRINTED 539
\weberpage{539}
of cyclic equations. This theorem can again be applied to the auxiliary equation of degree $s$, and one can continue in this way until this auxiliary equation reduces to the first degree. With this, then, one obtains the result:

\gesperrt{The solution of an Abelian equation can always be reduced to the solution of a sequence of cyclic equations whose degrees are divisors of the degree of the given equation.}

It is not necessary to include irreducibility in the definition of cyclic equations. We may therefore frame the definition generally as follows:

\gesperrt{An equation of degree $m$, $F(x) = 0$, with $m$ distinct roots is called a cyclic equation in the field $\Omega$ if its roots $\alpha, \alpha_1, \alpha_2 \ldots \alpha_{m-1}$ are not rational, but can be so arranged that the cyclic functions of the roots are rational in $\Omega$.}

If
\weq{7}{\pi = (\alpha, \alpha_1, \alpha_2 \ldots \alpha_{m-1})}
holds, then every function that admits the permutations of the cyclic group
\weq{8}{C = 1, \pi, \pi^2, \pi^3 \ldots \pi^{m-1}}
must be contained in $\Omega$.
The \gesperrt{Galois} group of a cyclic equation is either the period $C$ itself, and then the equation is irreducible, or it is a divisor of $C$; in that case, if $e$ and $f$ are two integer factors of $m$ and
\[
m = ef
\]
it consists of the permutations
\weq{9}{C_e = 1, \pi^e, \pi^{2e} \ldots \pi^{(f-1)e},}
and the cyclic equation splits into $e$ factors of degree $f$. For if we take a function $\psi(\alpha, \alpha_1 \ldots \alpha_{m-1})$ \gesperrt{belonging} to the group $C$, then by hypothesis it is equal to a quantity in $\Omega$, which we denote by $a$. No permutation not contained in $C$ is then applicable to the rational equation $\psi = a$, and consequently the group of the equation can contain no substitutions other than those which occur in $C$. Now if $\pi$ occurs in the group of the equation, then that group is identical with $C$, and since $C$ is transitive, the equation is irreducible. But if $\pi^e$ is the lowest power of $\pi$ which occurs in the
%CONT>>
